
% Ultra-minimal Overleaf project with NO BibTeX (thebibliography only)
\documentclass[11pt]{article}
\usepackage[T1]{fontenc}
\usepackage[utf8]{inputenc}
\usepackage{microtype}
\usepackage{geometry}
\geometry{margin=1in}
\usepackage{amsmath,amssymb,mathtools}
\usepackage{bm}
\usepackage{graphicx}
\usepackage{xcolor}
\usepackage[colorlinks=true,linkcolor=blue,citecolor=purple,urlcolor=teal]{hyperref}
\usepackage{enumitem}

% Soften line breaking to reduce overfull warnings
\emergencystretch=1em

% ===== Macros =====
\newcommand{\g}{\boldsymbol{g}}
\newcommand{\K}{\boldsymbol{K}}
\newcommand{\Top}{\mathsf{Top}}
\newcommand{\Stat}{\mathsf{Stat}}
\newcommand{\Dyn}{\mathsf{Dyn}}
\newcommand{\Q}{\mathcal{Q}}
\newcommand{\B}{\mathcal{B}}
\newcommand{\Comp}{\mathcal{C}\!omp}
\newcommand{\Konst}{\mathcal{K}\!onst}
\newcommand{\Fold}{\mathcal{F}\!old}
\newcommand{\Unfold}{\mathcal{U}\!nfold}

\title{Topographical Orthogonal Generative Theory (TOGT):\\A Computable Higher-Order Logic and Category-Theoretic Framework}
\author{Pablo Grossi}
\date{March 2026}

\begin{document}
\maketitle

\begin{abstract}
We present TOGT, a computable higher-order logic and category-theoretic framework unifying generative processes across domains. The core is a four-operator algebra (Compression, Constraint, Folding, Unfolding) organized in an algorithmically closed pipeline $S(L)\to\Phi\to\g\to\K(L)\to\text{Space}$. We formalize $\g=\nabla\Phi$, its sixfold composition $\g^{6}$ with a testable invariant (33 under constraints), and outline $\g^{64}$ (Kether Orthogon) as a 4-functor yielding a 4D crystalline spine.
\end{abstract}

\section{Background}
Complex systems exhibit directional, constrained, recursive generation. TOGT supplies a computable logic and category semantics to map patterns across scales.

\section{Framework}
Sorts: $\Top,\Stat,\Dyn$; potentials $\Phi$, flows $\g=\nabla\Phi$, curvature $\K$; operators $\{\Comp,\Konst,\Fold,\Unfold\}$. Aims include formal proofs, software, and validations across domains.

\section{Minimal Model}
$\Phi(x;\lambda)=\tfrac14 x^4-\tfrac12\lambda x^2$, $\dot{x}=\lambda x-x^3$.

\section{Notes}
See \cite{castelvecchi2025microsoft,ismail2022ewsph} for context; toolchains include TTK \cite{ttk}.

% ==== Manual bibliography (no BibTeX) ====
\begin{thebibliography}{9}
\bibitem{castelvecchi2025microsoft}
D. Castelvecchi. Microsoft claims quantum-computing breakthrough---but some physicists are sceptical. \emph{Nature}, 2025.

\bibitem{ismail2022ewsph}
M. M. S. Ismail et al. Early Warning Signals of Financial Crises Using Persistent Homology. \emph{Frontiers in Applied Mathematics and Statistics}, 2022.

\bibitem{ttk}
Topology ToolKit (TTK). Project site and repository, 2019--2026.
\end{thebibliography}

\end{document}
